\subsection{Optical Kerr Effect \& Self Steepening}
When the laser beam passes through the dielectric material, it induces a microscopic displacement of charges, forming oscillating electric dipoles that combine to produce macroscopic polarization. This results in a change in the refractive index. When the electric field is sufficiently intense, the response becomes nonlinear, giving rise to new propagation effects such as self-focusing and self-modulation, as well as plasma defocusing.
To determine the nonlinear polarization and the resulting Kerr effect within the SI system, we begin with the definition of the electric displacement field:
\begin{equation}
\mathbf{D} = \varepsilon_0 \mathbf{E} + \mathbf{P}
\end{equation}
where $\mathbf{P}$ can be expanded to include the third-order susceptibility contribution. Applied to a monochromatic electric field $\mathbf{E} = \mathcal{E}e^{i(k_0z-\omega_0t)} + \text{c.c.}$, this expression becomes:
\begin{equation}
\mathbf{P} = \varepsilon_0 \left( \chi^{(1)}\mathbf{E} + \chi^{(3)}\mathbf{E}^3 \right)
\end{equation}
By expanding $\mathbf{E}^3$ and neglecting the $3\omega_0$ terms associated with third-harmonic generation (THG), we isolate the term oscillating at the fundamental frequency, which represents the medium's nonlinear response at angular frequency $\omega_0$:
\begin{equation}
\mathbf{P} \approx \varepsilon_0 \left[ \chi^{(1)} + \frac{3}{4}\chi^{(3)}\vert{}\mathcal{E}\vert{}^2 \right] \mathbf{E}
\end{equation}
Substituting this expression back into the displacement field equation yields:
\begin{equation}
\mathbf{D} = \varepsilon_0 \left( 1 + \chi^{(1)} + \frac{3}{4}\chi^{(3)}\vert{}\mathcal{E}\vert{}^2 \right) \mathbf{E} = \varepsilon_0 \left( n_0^2 + \frac{3}{4}\chi^{(3)}\vert{}\mathcal{E}\vert{}^2 \right) \mathbf{E}= \varepsilon_0 \varepsilon_r\mathbf{E}
\end{equation}
where $n_0^2 = 1 + \chi^{(1)}$. The effective refractive index $n$ is defined as the square root of the relative permittivity: $n = \sqrt{\varepsilon_r}=\sqrt{n_0^2 + \frac{3}{4}\chi^{(3)}\vert{}\mathcal{E}\vert{}^2}$. For weak nonlinearities, a first-order Taylor expansion yields:
\begin{equation}
n \approx n_0 + \frac{3\chi^{(3)}}{8n_0}\vert{}\mathcal{E}\vert{}^2
\end{equation}
Comparing this expression with the standard form $n = n_0 + n_2 I$, where the intensity is $I = \frac{1}{2}n_0 c \varepsilon_0 \vert{}\mathcal{E}\vert{}^2$, we identify the nonlinear refractive index coefficient:
\begin{equation}
n_2 = \frac{3\chi^{(3)}}{4n_0^2 c \varepsilon_0}
\end{equation}
This refractive index therefore depends on the intensity. This is the origin of the optical Kerr effect. Since the intensity $I(r)$ of a Gaussian laser beam is higher at the center than at the periphery, the beam undergoes a higher refractive index at the center than at the edges, provided that $\chi^{(3)} > 0$. This spatial variation of the refractive index acts exactly like a gradient-index lens. The medium causes the wavefronts to converge toward the high-intensity axis. This phenomenon, known as self-focusing, is governed by the effective nonlinear polarization \cite{mao2004}
\begin{equation}
\mathbf{P}_{\text{NL}} = \frac{3}{4}\varepsilon_0 \chi^{(3)}\vert{}\mathcal{E}\vert{}^2 \mathbf{E}
\end{equation}
The phase of a femtosecond laser pulse in dielectric is dependant of the nonlinear refractive index $\Phi(t) = nk_0z-\omega_0 t = n_0k_0z +n2k_0zI -\omega_0 t$. Since the intensity of such impulsion is short enough the time evolution of the intensity has an effect on the frequency components of the pulse. The phase evolution undergoes a self steepening effect : The frequency evolution is shifted as the pulse propagate into the material. If the intensity profile on the input is symetric, that would be the case of a perfect gaussian beam, the spectrum can therfore remain symetric.

\begin{equation}
\frac{d\Phi}{dt}= -\omega_0+n_2z\frac{\omega_0}{c}\frac{dI(t)}{dt}
\end{equation}
